% @author Lluis Gesa (lluis.gesa@uab.cat)
%
% Enginyeria Del Software
% -----------------------------------------------------------------------------------



% definim quin estil general volem per el document
\documentclass[11pt]{article}

% indiquem quins paquets amb funcionalitat extra volem fer servir
\usepackage{a4wide, amsmath, tabularx, colortbl, fancyhdr, graphicx, lastpage, caption}
\usepackage[utf8]{inputenc}

% Indiquem que comencem a definir el document i els seus continguts.
\begin{document}

% Definim valors i propietats, creem macros que ens ajudaran a
% agilitzar coses repetides:

\newcommand{\NomProjecte}{Perdidos en la ignorancia}
\newcommand{\NomDocument}{Acta Reunió}
\newcommand{\IdentificadorDocument}{GG-SL-N-MoM-01}
\newcommand{\VersioDocument}{v1.0}

% Definim una macro que cada cop que insertem \requeriment, creara una
% taula amb el contingut passat per paràmetre
\newcommand{\requeriment}[6]{
\begin{center}
	\begin{tabularx}{\linewidth}{|p{2.5cm}|X|}
	\hline
	\cellcolor[gray]{.8} \textbf{ID} & \cellcolor[gray]{.9} #1\\
	\hline
	\cellcolor[gray]{.8} \textbf{Title} & #2 \\
	\hline
	\cellcolor[gray]{.8} \textbf{Description} & #3\\
	\hline
	\cellcolor[gray]{.8} \textbf{Priority} & #4 \\
	\hline
	\cellcolor[gray]{.8} \textbf{Verification} & #5 \\
	\hline
	\cellcolor[gray]{.8} \textbf{Parents} & #6\\
	\hline
	\end{tabularx}
\end{center}
}

% Fent servir propietas del paquet fancyhdr per crear una capçalera
% que es mostrarà cada pàgina
\pagestyle{fancy}
\fancyhead[LO,LE]{\includegraphics[width = 3.5cm]{./imatges/uab-logo.png}}
\fancyhead[CO,CE]{\bf \NomProjecte \\ \vspace{\baselineskip}
  \NomDocument }
\fancyhead[RO,RE]{
{\bf Ref : \IdentificadorDocument} \\
Versió: \VersioDocument\\
Data: \today \\
Pag. \thepage ~of \pageref{LastPage}
}
\fancyfoot[LO,LE,CO,CE,RO,RE]{}
\textheight=19cm


%
% Creem la pàgina inicial :

% Insertem el nom del document en una taula centrada feng servir una
% font gran \huge

\begin{center}
\begin{tabular*}{15cm}{p{14.5cm}}
\\
\\
\\
\rowcolor[gray]{.9}\begin{center}\huge{\bf{\NomDocument}}\end{center}\\
\\
\end{tabular*}
\end{center}


%
% Insertem una taula que contindrà les abreviacions que hi hagi al
% nostre document
%
\small{%
\begin{center}
\begin{tabular}{|p{2cm}|p{4.75cm}|p{2cm}|p{4.75cm}|}
\hline
\multicolumn{4}{|>{\columncolor[gray]{.9}[2mm][2mm]}c|}{{\bf Abreviatures}} \\
\hline
UAB & Universitat Autonoma de Barcelona & ES & Enginyeria del Software \\
\hline
& & & \\
\hline
\end{tabular}
\end{center}
}

%
% Insertem una taula amb el historial de canvis del document
%
\small{%
\begin{center}
\begin{tabular}{|p{2cm}|p{2cm}|p{7.5cm}|p{2cm}|}
\hline
\multicolumn{4}{|>{\columncolor[gray]{.9}[2mm][2mm]}c|}{{\bf Historial de Revisions}} \\
\hline
{\bf Version} & {\bf Date} & {\bf Comments} & {\bf Autor} \\
\hline
1.0 & 08-Abril-2019 & Sprint 1 &
G43-2-02 \\
\hline
\end{tabular}
\end{center}
}




% Insertem en una nova pàgina, els indexs : de continguts, de figures
% i de taules, les comandes \tableofcontents ,  \listoffigures,
% \listoftables són gestionades per els paquets fets servir, i per
% defecte escriuen els titols de la secció en angles, per canviar-los,
% hem de 'renovar les comandes' :

\renewcommand{\contentsname}{Índex de continguts}
\renewcommand{\listfigurename}{Índex de figures}
\renewcommand{\listtablename}{Índex de taules}

% Ara li diem al LaTeX que inserti la informació
\newpage
\tableofcontents
\listoffigures
\listoftables

% Insertem les seccions que volem
\newpage

\section{Participants}

\begin{center}
\begin{tabular}{|c|c|}
\hline
{\cellcolor[gray]{.8} \bf Nom} & {\cellcolor[gray]{.8} \bf Cognom}  \\
\hline
Daniel & Montesinos \\
\hline

Aleksandra & Jovanovic \\
\hline

Núria & Navarro \\
\hline

Adrián & Moreno \\
\hline

Desirée & Pasión \\
\hline
\end{tabular}
\captionof{table}{Participants}
\end{center}

\section{Agenda}

\begin{itemize}
\subsection{Sprint 1}
\item Bienvenida: Inicio de la sesión de presentaciones.
\item Presentación: Exposición de cada grupo de la tarea realizada durante el Sprint 1.
\item Organización: Organización de las tareas a realizar durante el Sprint 2.
\begin{enumerate}
\item Documento MoM: Creación del documento
\item Acabar GameDoc: Faltaban algunas partes del documento
\item Presentación Sprint 2 : Creación de la presentación
\item Buscar información: Información referente a Bitbucket y su funcionamiento
\item Trazabilidad: Añadirla al documento correspondiente
\item Diagrama Casos de uso : Terminar el diagrama
\item Diagrama de clase: Crear el diagrama de clases
\item Correcciones documentos: Corregir versionado, nombres, etc.
\item Reasignar tipo de Prioridad y Verification Parents : Reasignar en documento de requerimientos
\item Crear y documentar documento UML clases: Crearlo mediante una plantilla anterior
\item Ampliar y terminar requerimientos más específicos
\end{enumerate}



\subsection{Sprint 2}
\item Presentaciones: Exposición de cada grupo de la tarea realizada durante el Sprint 2.
\item Organización: Organización de las tareas a realizar durante el Sprint 3.
\begin{enumerate}
	\item Documentar MoM v2.0: Documentar acta de la reunión.
	\item Diagrama de secuencia: Crear el diagrama de secuencia.
	\item Documentar diagrama de secuencia: Documentación del diagrama de secuencia.
	\item Diagrama de estados: Crear el diagrama de estados.
	\item Documentar diagrama de estados: Documentación del diagrama de estados.
	\item Diagrama de clase: Añadir cardinalidad al diagrama de clase.
	\item Documento GameDoc: Arreglar la versión.
	\item Documento GameDoc: Hacer apartado 2.
	\item Documento Requerimientos: Comprovar trazabilidad.
	\item Documento Requerimientos: Asignación correcta de Test o Diseño.
	\item Documento Requerimientos: Refinar los requerimientos.
	\item Presentación: Crear la presentación del Sprint 3.
	\item Revisar: Revisión de los documentos creados i/o modificados durante el Sprint (Especificaciones, Clases, GameDoc, Requerimientos, Secuencia, Estados)
	\item Hacer commits de las modificaciones en bitbucket.
\end{enumerate}







\end{itemize}

\section{Comentaris}


\subsection{Sprint 1}

En este documento se presentan las tareas a realizar en el Sprint 2 del proyecto. En este Sprint se planea terminar las tareas del Sprint 1 que quedaron pendientes y ampliar con diagramas de clases. Para esto se crearán nuevos documentos y se editarán los anteriores.


\subsection{Sprint 2}
En este documento se presentan las tareas a realizar en el Sprint 3 del proyecto. En este Sprint se planea terminar las tareas del Sprint 2 y del Sprint 1 que no eran correctas. También tiene como objetivo la creación de los diagramas de secuencia y de estado.


\subsection{Sprint 3}
En este documento se presentan las tareas a realizar en el Sprint 3 del proyecto. En este Sprint se planea terminar las tareas del Sprint 3 que no se terminaron y modificar aquellas que no eran correctas, como por ejemplo modificar algunos errores del diagrama de clases, de secuencia, y mejorar el diagrama de estados. También tiene como objetivo la creación de los diagramas de test.

\begin{center}
\begin{tabular}{|c|c|}
\hline
{\cellcolor[gray]{.8} \bf Acció Num} & {\cellcolor[gray]{.8} \bf Descripció}  \\
\hline
ACT-01 & La presentación ha ido de forma correcta. Se han descubierto nuevas necesidades y errores. \\
\hline
ACT-02 & La presentación ha ido de forma correcta. Hemos conocido nuevos errores. \\
\hline
\end{tabular}
\captionof{table}{Accions}
\end{center}


% Finalitzem el document.
\end{document}
